\chapter{A joint moment identity for one isotropic law}\label{ch-moment}

Chapters~\ref{ch-variational} and~\ref{ch-support} reduced the structural
problem to a single probability law. This chapter and the next bound it. From
here to the end of Proposition~\ref{final-certificate} the object of study is
an arbitrary Borel probability law on $\R^3$ with $\E XX^T=I_3$; the
maximising graphon enters only at the very end, in
Section~\ref{sec-conclusion}, where its canonical law is shown to satisfy the
extra radial hypothesis that the argument needs.

The two quantities to be controlled are
\[
 A=\E|X\cdot X'|,\qquad B=|\E X|^2,
\]
and the target is $A+B\le2c$. The classical route bounds them separately:
$A\le\gamma(3)=(1+\sqrt5)/2$ and $B\le1$, giving $A+B\le(3+\sqrt5)/2$ and
therefore the coefficient $(3+\sqrt5)/4$. That route is exhausted, because
both individual bounds are sharp and are attained by different laws. Any
improvement must exploit the fact that the two quantities cannot be large
simultaneously.

The device used here is a single exact identity rather than a chain of
estimates. Section~\ref{sec-moments} fixes the notation and decomposes the
law into a radial part $R=|X|$, a directional part $U=X/R$, the mean
$\mathbf m=\E X$, the second-moment tensor $M=\E[RUU^T]=\tfrac C3I_3+Z$, and
the regression vector $\xi=U-MX$. It records the scalar budget
$D=e+\alpha V+\beta J+\rho K^2+S$ coming from the classical absolute-value
majorant $h_0$, and proves a regression estimate that bounds the weighted
linear moment $L=\E[R(t-q_0)^2]$ in terms of $e,V,J$ with an explicit
nonnegative remainder $r_{\rm reg}$.

Section~\ref{sec-quartic} replaces $h_0$ by a signed quartic. The polynomial
$f$ is the degree-four Hermite interpolant with double contact to $|w|+w$ at
$w=-3/8$ and $w=1/2$ and a contact at $w=1$; its slack $\sigma=f-|w|-w$ is
nonnegative on $[-1,1]$ and, crucially, dominates the explicit quadratic
$(1-9w^2)/25$. Integrating the expansion of $f$ against the harmonic
decompositions produces the single exact identity~\eqref{budget7},
\[
 D=\delta_{\mathrm q}+b_ee+b_VV+b_JJ+a_3H^2+a_4K^2+P,
\]
in which the mean deficit $e$, the radial variance $V$, the regression error
$J$, the two harmonic moments $H,K$ and the slack moment $P$ all appear at
once. Lemma~\ref{radialpayment} then extracts from $P$ the linear payment
$\tfrac3{50}(V-9J)$ plus a nonnegative remainder $r_{\rm rad}$, under a
bounded-radius hypothesis with endpoint ratio at most $3$.

Nothing in this chapter is optimised yet; the identity is exact and every
remainder is retained. The next chapter pays for the coefficient.

\section{Moments and regression}\label{sec-moments}
% Ported verbatim from paper/top3_joint_moment_bound.tex lines 718-816.
% Only cross-reference wording was adapted; no mathematics was altered.
We now turn to the numerical estimate. Throughout this section the law
is arbitrary and isotropic. Bounded positive radii will be imposed only
when a radial comparison requires them. For later bookkeeping define
\begin{equation}\label{constants}
 d_0=-\frac{25624284539519}{53655078400000}
                  +\frac{\sqrt5}{2}-\frac{41\sqrt3}{120},\qquad
 d_1=d_0+\frac{103}{9000000},\qquad
 d=d_1+\frac1{3000000}.
\end{equation}
Then $d=(3+\sqrt5)/2-2c$, with $c$ as in~\eqref{main-constant}.
The intermediate constants separate the joint moment estimate from
the final atomic term.
Fix $\E XX^T=I_3$, let $R=|X|$, and put $U=X/R$ when $R>0$.
At $R=0$ choose any fixed unit vector for $U$. Define
\[
 \begin{gathered}
 C=\E R,\quad V=3-C^2,\quad \mathbf m=\E X,\quad B=|\mathbf m|^2,\quad e=1-B,\\
 A=\E|X\cdot X'|,\qquad D=\frac{3+\sqrt5}{2}-A-B.
 \end{gathered}
\]
Here $V,e\ge0$: $\E R^2=3$ and Cauchy--Schwarz gives
$(v\cdot \mathbf m)^2\le\E(v\cdot X)^2=1$ for every unit $v$.
All pair expectations below are over independent copies of this same law.
Set
\[
 M=\E[R UU^T]=\frac C3I_3+Z,\qquad
 J=1-\tr(M^2)=\frac V3-\|Z\|_F^2.
\]
The regression vector $\xi=U-MX$ satisfies
\begin{equation}\label{regression}
 \E[X\xi^T]=0,\quad \E\|\xi\|^2=J\ge0,\quad
 \E UU^T=M^2+\E[\xi\xi^T].
\end{equation}
In particular $J$ is an actual regression error, not a freely chosen moment.

Here are the harmonic normalizations used throughout. All tensor inner
products sum over all ordered indices. Put
\[
 \Omega_{ijkl}=\frac{\delta_{ij}\delta_{kl}+\delta_{ik}\delta_{jl}
                           +\delta_{il}\delta_{jk}}{15},
 \qquad
 ({\cal L}_3(y))_{ijk}
       =\frac{y_i\delta_{jk}+y_j\delta_{ik}+y_k\delta_{ij}}5.
\]
For symmetric traceless $Z$, put
\[
 ({\cal L}_4(Z))_{ijkl}=\frac17(\delta_{ij}Z_{kl}+\delta_{ik}Z_{jl}
 +\delta_{il}Z_{jk}+\delta_{jk}Z_{il}+\delta_{jl}Z_{ik}+\delta_{kl}Z_{ij}).
\]
For unit $u$, set
\[
 \begin{gathered}
 \Psi_3(u)=u^{\otimes3}-{\cal L}_3(u),\qquad
 \Psi_4(u)=u^{\otimes4}-\Omega-{\cal L}_4(uu^T-I_3/3),\\
 T_i=\E[R\Psi_i(U)],\qquad H=\|T_3\|_F,\qquad K=\|T_4\|_F.
 \end{gathered}
\]
Contraction in two indices gives the orthogonal decompositions
\[
 \E[R U^{\otimes3}]=T_3+{\cal L}_3(\mathbf m),\qquad
 \E[R U^{\otimes4}]=T_4+C\Omega+{\cal L}_4(Z).
\]
Indeed $\operatorname{ct}\Omega=I_3/3$, $\|\Omega\|^2=1/5$,
$\operatorname{ct}{\cal L}_4(Z)=Z$, and
$\langle T,{\cal L}_4(Z)\rangle=(6/7)\langle\operatorname{ct}T,Z\rangle$.
Also $\operatorname{ct}{\cal L}_3(y)=y$ and
$\langle N,{\cal L}_3(y)\rangle=(3/5)\operatorname{ct}N\cdot y$.
It follows directly that
\begin{align}\label{harmonics}
 \langle\Psi_3(u),\Psi_3(v)\rangle&=g_3(u\cdot v),&
 g_3(t)&=t^3-3t/5,&\|\Psi_3(v)\|^2&=2/5,\notag\\
 \langle\Psi_4(u),\Psi_4(v)\rangle&=g_4(u\cdot v),&
 g_4(t)&=t^4-6t^2/7+3/35,&\|\Psi_4(v)\|^2&=8/35.
\end{align}
The tensor moments involve $R$, not $R^3$ or $R^4$, and are integrable
under $\E R^2=3$.

Write $s=\sqrt5$ and define the scalar absolute-value majorant and its slack by
\[
 a=\frac{1+s}{6},\quad b=\frac{7s-5}{8},\quad\rho=\frac{15s-25}{16},
 \qquad h_0(z)=\frac{(1-z)(sz-1)^2(sz+s+2)}{2(s+1)^2}.
\]
Expansion gives, for $|t|\le1$,
\[
 a+b(t^2-1/3)-\rho(t^4-1/5)-|t|=h_0(|t|)\ge0.
\]
This majorant is from~\cite[Lemma 4.4]{STW}.
Let $S=\E[RR'h_0(|U\cdot U'|)]$. The decompositions above give
\[
 A=aC^2+b\|Z\|^2-\rho\left(K^2+\frac67\|Z\|^2\right)-S.
\]
Consequently
\begin{equation}\label{budget4}
 D=e+\alpha V+\beta J+\rho K^2+S,
 \qquad \alpha=\frac{2s-1}{14},\quad\beta=\frac{s+10}{14}.
\end{equation}
Useful rational bounds are $\alpha\ge26/105$, $\beta\ge6/7$,
and $\rho\ge1/2$, using $s>67/30$, $s\ge2$, and $s\ge11/5$.


\subsection{A regression estimate}
% Ported verbatim from paper/top3_joint_moment_bound.tex lines 818-852.
% Only cross-reference wording was adapted; no mathematics was altered.
Suppose $B>0$. Put
\[
 \begin{gathered}
 x=\sqrt B,\quad v=\mathbf m/x,\quad t=U\cdot v,\quad Y=Rt,\\
 q_0=1/\sqrt3,\quad q=C/(3x),\quad L=\E[R(t-q_0)^2].
 \end{gathered}
\]
Then $\E Y=x$, $\E Y^2=1$, $\operatorname{Var}Y=e$.
If $w_0=\E\xi$, \eqref{regression} gives
$w_0=\E[(1-\mathbf m\cdot X)\xi]$, hence $|w_0|^2\le eJ$.
For traceless symmetric $3$ by $3$ matrices,
$|Zv|^2\le(2/3)\|Z\|^2$; the other two eigenvalues accompanying an
eigenvalue $z$ have squared sum at least $z^2/2$.
Writing $J_v=\E(v\cdot\xi)^2\le J$, we obtain
\begin{align*}
 \E(t-q)^2
 &=|Mv-q\mathbf m|^2+q^2e+J_v-2qv\cdot w_0\\
 &\le 2V/9+J/3+q^2e+2q\sqrt{eJ}=:F_0,\\
 \E(Y-qR)^2&\le V/3+3q^2e+2q\sqrt{eV}=:G_0.
\end{align*}
Here $Mv-q\mathbf m=Zv$ and $\E RY\ge Cx-\sqrt{eV}$.
Cauchy--Schwarz gives $L_q:=\E[R(t-q)^2]\le\sqrt{F_0G_0}$.
Note that $Rt=R(U\cdot v)=(RU)\cdot v=X\cdot v=Y$, so $\E[Rt]=\E Y=x$;
this is the only place at which the regression vector enters the exact
identity below. With $\zeta=C/\sqrt3$,
\[
 L_q-L=q_0(\zeta-x)^2(\zeta+2x)/x^2\ge0.
\]
Thus, writing $u=\sqrt e$, $z=\sqrt V$, $j=\sqrt J$,
\begin{equation}\label{Lbar}
 L\le\overline L:=\frac34F_0+\frac13G_0
 =\frac{7q^2}{4}u^2+\frac5{18}z^2+\frac14j^2
                       +\frac{2q}{3}uz+\frac{3q}{2}uj.
\end{equation}
Set $r_{\rm reg}=\overline L-L\ge0$. This remainder will be kept
in the joint identity. No lower radius bound has been used here.


\section{A signed quartic and a joint defect identity}\label{sec-quartic}
% Ported verbatim from paper/top3_joint_moment_bound.tex lines 854-955.
% Only cross-reference wording was adapted; no mathematics was altered.
Define
\[
 f(w)=\frac{8352+43806w+54016w^2-12928w^3-10240w^4}{41503},
 \qquad \sigma(w)=f(w)-|w|-w.
\]
The polynomial is the degree-at-most-four Hermite interpolant with double
contacts to $|w|+w$ at $-3/8$ and $1/2$, and a contact at $1$.
Uniqueness follows because a difference of two such polynomials would have
five zeros counted with multiplicity. The following explicit identities prove
the needed inequalities; interpolation alone is not a majorant proof.
\begin{lemma}\label{scalar}
For every $w\in[-1,1]$,
\begin{equation}\label{sigmapoint}
 \sigma(w)\ge0,\qquad \sigma(w)\ge\frac{1-9w^2}{25}.
\end{equation}
\end{lemma}
\begin{proof}
Direct expansion gives
\[
 41503\sigma(w)=
 \begin{cases}
 (1-w)(2w-1)^2(8352+2560w),&0\le w\le1,\\
 (8w+3)^2(928-82w-160w^2),&-1\le w\le0.
 \end{cases}
\]
The final factor on the negative interval is concave and at least $850$.
This proves nonnegativity, including the join at $0$ and both endpoints.
The second assertion follows immediately when $|w|\ge1/3$.
For $w=-z$, $0\le z\le1/3$, write
\[
 \sigma(-z)=(z-3/8)^2
          \frac{59392+5248z-10240z^2}{41503}
       \ge\frac{10}{7}(z-3/8)^2.
\]
Indeed $z(5248-10240z)\ge0$ on this interval and
$59392/41503>10/7$. The remaining comparison is the exact square
\[
 \frac{10}{7}(z-3/8)^2-\frac{1-9z^2}{25}
  =\frac{313}{175}\left(z-\frac{375}{1252}\right)^2
                         +\frac{109}{250400}>0.
\]
For $w=z$, $0\le z\le1/3$, the positive-interval factorization gives
\[
 \sigma(z)=(z-1/2)^2\frac{(1-z)(33408+10240z)}{41503}
              \ge\frac12(z-1/2)^2,
\]
because the coefficient is at least $22272/41503>1/2$. Now use
\[
 \frac12(z-1/2)^2-\frac{1-9z^2}{25}
   =\frac{43}{50}\left(z-\frac{25}{86}\right)^2+\frac{53}{4300}>0.
\]
These intervals exhaust $[-1,1]$.
\end{proof}

Put $P=\E[RR'\sigma(U\cdot U')]$ for the same actual law as before, and set
\[
 \begin{gathered}
 a_3=\frac{12928}{41503},\quad a_4=\frac{10240}{41503},\quad
 b_e=\frac{16386}{18865},\\
 b_V=\frac{64608}{290521},\quad b_J=\frac{316672}{290521},\quad
 \delta_{\mathrm q}=\frac{3+\sqrt5}{2}-\frac{544886}{207515}.
 \end{gathered}
\]
Integration of the expanded $f$, using \eqref{harmonics}, gives the single
exact joint budget
\begin{equation}\label{budget7}
 D=\delta_{\mathrm q}+b_e e+b_V V+b_J J+a_3H^2+a_4K^2+P.
\end{equation}
For example, writing $f=c_0+c_1w+c_2w^2-a_3w^3-a_4w^4$,
the coefficients are $b_e=c_1-3a_3/5$, $b_V=c_0+3a_4/35$,
$b_J=c_2-6a_4/7$, and the integrated constant is
$3(c_0-a_4/5)+b_e+c_2=544886/207515$.
All tensors are $R$-weighted and $\sigma$ is bounded, so this equality
uses no radial moment beyond the second one. Identity~\eqref{budget4} will only justify a small-$D$ branch; it is
not added as an independent resource to~\eqref{budget7}.

\begin{lemma}\label{radialpayment}
Suppose $0<r_-\le R\le r_+<\infty$ almost surely and $r_+/r_-\le3$.
Then
\begin{equation}\label{payment}
 P=\frac3{50}(V-9J)+r_{\rm rad},\qquad r_{\rm rad}\ge0.
\end{equation}
\end{lemma}
\begin{proof}
For $w=U\cdot U'$, the exact pair identities from the full Gram are
\[
 V=\frac12\E(R-R')^2,\qquad
 J=\frac12\E[(R-R')^2w^2].
\]
The latter uses $\E[R^2w^2]=\E[(R')^2w^2]=1$ and
$\E[RR'w^2]=\tr(M^2)$; it assumes no independence of radius and direction.
The ratio bound gives $(R-R')^2\le(4/3)RR'$. Use the two conclusions
of Lemma~\ref{scalar} in this order:
\[
 P\ge\frac34\E[(R-R')^2\sigma(w)]
   \ge\frac3{100}\E[(R-R')^2(1-9w^2)]
    =\frac3{50}(V-9J).
\]
The expression $V-9J$ need not be nonnegative. Define $r_{\rm rad}$
as the actual difference; only this difference is claimed nonnegative.
\end{proof}

